\section*{Hoofdstuk \ref{ch:b1:isostasy-gravity} --- Isostasie en zwaartekrachtsanomalieën}

\begin{solution}{exo:b1:isostasy-gravity:1}
(a) $\qty{-150}{\milli\gal} = -150\times\qty{e-5}{m.s^{-2}} = \qty{-1.5e-3}{m.s^{-2}}$,
en $\num{1.5e-3}/9.81 = \num{1.53e-4}$ van $g$. (b) $\qty{1}{\milli\gal} =
\qty{e-5}{m.s^{-2}}$, \num{1.02e-6} van $g$; $\qty{1}{\micro\gal} =
\qty{e-8}{m.s^{-2}}$, \num{1.02e-9} van $g$: een zwaartekrachtmeter die een microgal onderscheidt, meet de zwaartekracht tot op ongeveer één deel per miljard.
\end{solution}

\begin{solution}{exo:b1:isostasy-gravity:2}
(a) Aan de evenaar verdwijnen beide sinussen: $\gamma = \qty{9.780327}{m.s^{-2}}$. Op
\ang{45} is $\sin^2\varphi = 0.5$ en $\sin^2 2\varphi = 1$:
$\gamma = 9.780327\times(1 + 0.0026512 - 0.0000058) = \qty{9.806200}{m.s^{-2}}$.
Aan de polen is $\sin^2\varphi = 1$ en $\sin 2\varphi = 0$:
$\gamma = 9.780327\times1.0053024 = \qty{9.832186}{m.s^{-2}}$.
(b) $9.832186 - 9.780327 = \qty{0.051859}{m.s^{-2}}$, dat is
\qty{5186}{\milli\gal}. (c) $3392/5186 = 0.65$: het centrifugale effect verklaart
ongeveer twee derde van het verschil; de afplatting, die de polen dichter bij het middelpunt brengt, de rest.
\end{solution}

\begin{solution}{exo:b1:isostasy-gravity:3}
$r = h\rho_c/(\rho_m - \rho_c) = 3\times2800/500 = \qty{16.8}{km}$. De \omterm{def:b1:layered-earth:layers}{korst} is
dikker met $h + r = \qty{19.8}{km}$ (dat is $h\rho_m/(\rho_m - \rho_c) =
3\times3300/500$).
\end{solution}

\begin{solution}{exo:b1:isostasy-gravity:4}
(a) $1/0.3086 = \qty{3.24}{m}$. (b) $0.01/0.3086 = \qty{0.032}{m}$: de hoogten
moeten bekend zijn tot ongeveer \qty{3}{cm}, en daarom registreert elke meetcampagne de hoogte van elk station met satellietpositionering of waterpassing.
\end{solution}

\begin{solution}{exo:b1:isostasy-gravity:5}
(a) $\rho = 2700\times100/(100 + h)$: \qty{2673}{kg.m^{-3}} voor \qty{1}{km},
\qty{2621}{kg.m^{-3}} voor \qty{3}{km}, \qty{2571}{kg.m^{-3}} voor \qty{5}{km}.
(b) $\rho = (2700\times100 - 1030\times4)/96 = \num{265880}/96 =
\qty{2770}{kg.m^{-3}}$: de kolom onder de zee is dichter. (c) Het model van Pratt vraagt om laterale veranderingen van dichtheid bij een constante diepte. Onder de oceaanbodem bestaan zulke veranderingen: jonge, hete \omterm{def:b1:layered-earth:lithosphere}{lithosfeer} nabij een rug is lichter dan oude, koude \omterm{def:b1:layered-earth:lithosphere}{lithosfeer} ver ervandaan (\cref{ch:b1:heat-budget}). Onder de continenten daarentegen blijkt de dikte van de \omterm{def:b1:layered-earth:layers}{korst} met het reliëf te veranderen.
\end{solution}

\begin{solution}{exo:b1:isostasy-gravity:6}
(a) Station 1: $g - \gamma = 9.807751 - 9.808911 = \qty{-0.001160}{m.s^{-2}} =
\qty{-116.0}{\milli\gal}$; \omterm{def:b1:isostasy-gravity:anomalies}{vrije-luchtcorrectie} $0.3086\times350 =
\qty{108.0}{\milli\gal}$, dus $\Delta g_{\mathrm{FA}} = \qty{-8.0}{\milli\gal}$;
\omterm{def:b1:isostasy-gravity:anomalies}{Bouguer-correctie} $0.112\times350 = \qty{39.2}{\milli\gal}$, dus
$\Delta g_{\mathrm B} = \qty{-47.2}{\milli\gal}$.
Station 2: $g - \gamma = \qty{-535.9}{\milli\gal}$; \omterm{def:b1:isostasy-gravity:anomalies}{vrije-luchtcorrectie}
$0.3086\times1850 = \qty{570.9}{\milli\gal}$, dus $\Delta g_{\mathrm{FA}} =
\qty{+35.0}{\milli\gal}$; \omterm{def:b1:isostasy-gravity:anomalies}{Bouguer-correctie} $0.112\times1850 =
\qty{207.2}{\milli\gal}$, dus $\Delta g_{\mathrm B} = \qty{-172.2}{\milli\gal}$.
(b) Op station 2 is de \omterm{def:b1:isostasy-gravity:anomalies}{Bouguer-anomalie} sterk negatief: onder de bergen is er een massatekort, een wortel van lichte \omterm{def:b1:layered-earth:layers}{korst}. De \omterm{def:b1:isostasy-gravity:anomalies}{vrije-luchtanomalie}, klein vergeleken met de \omterm{def:b1:isostasy-gravity:anomalies}{Bouguer-correctie}, toont dat het tekort de bergen bijna in evenwicht houdt: zij zijn grotendeels gecompenseerd (tot ongeveer $172/207$, of \qty{83}{\percent}, als het plaatmodel voor de hele keten geldt).
\end{solution}

\begin{solution}{exo:b1:isostasy-gravity:7}
(a) Zoals in het bewijs van de propositie levert een dunne laag op de afstand $z$
een bijdrage $G\rho\,2\pi z\,h\int_0^a s\,(s^2 + z^2)^{-3/2}\,\mathrm ds$, nu met de integraal gestopt bij $s = a$:
$\int_0^a s\,(s^2 + z^2)^{-3/2}\,\mathrm ds = 1/z - 1/\sqrt{z^2 + a^2}$. Dus
$g = 2\pi G\rho h\,\bigl(1 - z/\sqrt{z^2 + a^2}\bigr)$, wat naar de plaatwaarde gaat voor $a \to \infty$. (b) $z/\sqrt{z^2 + a^2} = 0.1$ geeft
$\sqrt{z^2 + a^2} = 10z$, $a = \sqrt{99}\,z = \qty{9.95}{km}$ voor
$z = \qty{1}{km}$. (c) De rots binnen ongeveer tien keer haar diepte onder het station levert \qty{90}{\percent} van de aantrekking: de reliëfcorrectie heeft het reliëf nodig over een straal van enkele kilometers tot tientallen kilometers rond elk station, niet alleen eronder.
\end{solution}

\begin{solution}{exo:b1:isostasy-gravity:8}
(a) $r' = d(\rho_c - \rho_w)/(\rho_m - \rho_c) = 4\times1770/500 =
\qty{14.16}{km}$, en de \omterm{def:b1:layered-earth:layers}{korst} is $35 - 4 - 14.16 = \qty{16.84}{km}$ dik.
(b) Het model van Airy met één korstdichtheid voorspelt een oceanische \omterm{def:b1:layered-earth:layers}{korst} die meer dan drie keer te dik is. De oceanische \omterm{def:b1:layered-earth:layers}{korst} is geen uitgedunde continentale \omterm{def:b1:layered-earth:layers}{korst}: zij is een ander, dichter gesteente, en de twee niveaus van de aarde zijn twee soorten \omterm{def:b1:layered-earth:layers}{korst}, zoals de \omterm{def:b1:isostasy-gravity:hypsometry}{hypsometrische kromme} doet vermoeden.
\end{solution}

\begin{solution}{exo:b1:isostasy-gravity:9}
(a) Eén boogseconde is $\pi/(180\times3600) = 1/\num{206265}$ van een radiaal:
\qty{5.236}{''} $= \qty{2.538e-5}{rad}$ en \qty{15.885}{''} $=
\qty{7.701e-5}{rad}$. (b) $a_h = g\theta$: $9.79\times\num{2.538e-5} =
\qty{2.485e-4}{m.s^{-2}} = \qty{24.9}{\milli\gal}$ waargenomen, en
$9.79\times\num{7.701e-5} = \qty{75.4}{\milli\gal}$ berekend uit de zichtbare
bergen. (c) $1 - 5.236/15.885 = 0.670$: de compensatie heeft ongeveer twee derde van de aantrekking van de zichtbare bergen opgeheven.
\end{solution}

\begin{solution}{exo:b1:isostasy-gravity:10}
(a) Buiten de bol trekt een bol aan als een puntmassa in haar middelpunt, met
$GM/(x^2 + z^2)$ langs de lijn naar het middelpunt; de verticale component is dat maal de cosinus $z/\sqrt{x^2 + z^2}$, wat $GMz/(x^2 + z^2)^{3/2}$ geeft.
(b) $M = \tfrac43\pi(\qty{2000}{m})^3\times(\qty{-200}{kg.m^{-3}}) =
\qty{-6.70e12}{kg}$, en boven het middelpunt $GM/z^2 = \num{6.674e-11}\times
(\num{-6.70e12})/(4000)^2 = \qty{-2.80e-5}{m.s^{-2}} = \qty{-2.80}{\milli\gal}$.
(c) $(1 + x^2/z^2)^{-3/2} = 1/2$ geeft $x = z\sqrt{2^{2/3} - 1} = 0.766\,z =
\qty{3.07}{km}$. (d) De halve breedte op halve hoogte van het profiel geeft de diepte: $z = 1.305\,x_{1/2}$. Een dieper lichaam geeft een bredere, vlakkere anomalie.
\end{solution}

\begin{solution}{exo:b1:isostasy-gravity:11}
(a) $V = \tfrac13\pi R^2 H = \tfrac13\pi(\qty{15e3}{m})^2(\qty{3e3}{m}) =
\qty{7.07e11}{m^3}$, en $\Delta m = (2800 - 1030)V = \qty{1.25e15}{kg}$.
(b) Het zwaartepunt ligt \qty{0.75}{km} boven de vloer, op $d = 5 - 0.75 = \qty{4.25}{km}$ onder het zeeoppervlak. Daar is de potentiaalanomalie $T = -G\Delta m/d$, dus $N = G\Delta m/(\gamma d) =
\num{6.674e-11}\times\num{1.25e15}/(9.81\times4250) = \qty{2.0}{m}$.
(c) Het grootste deel van de massa van de kegel ligt in zijn brede onderste deel, op \qty{5}{km} tot \qty{16}{km} van het punt boven de top: het gemiddelde van $1/r$ over de massa is kleiner dan $1/\qty{4.25}{km}$, en de puntmassa overschat de potentiaal. Door dunne schijven op te tellen komt men op ongeveer \qty{1.1}{m}.
(d) Een wortel van lichte \omterm{def:b1:layered-earth:layers}{korst} eronder zou een tekort van dezelfde massa dieper leggen; de twee potentialen heffen elkaar bijna op, en de bult zou veel kleiner zijn, maar niet nul, omdat het tekort verder van het oppervlak ligt dan het overschot.
\end{solution}

\begin{solution}{exo:b1:isostasy-gravity:12}
(a) De opheffingssnelheid is $-\mathrm dw/\mathrm dt = w_0\,e^{-t/\tau}/\tau = w/\tau$.
(b) $w = 818\,e^{-10000/4000} = 818\times0.0821 = \qty{67.1}{m}$, en de snelheid is
$\qty{67.1}{m}/\num{4000}\,\unit{yr} = \qty{16.8}{mm.yr^{-1}}$, meer dan anderhalf keer de gemeten waarde. (c) Met $\tau = 3100$ jaar is $w =
818\,e^{-3.226} = \qty{32.5}{m}$ en de snelheid \qty{10.5}{mm.yr^{-1}}; met
$\tau = \num{68000}$ jaar is $w = 818\,e^{-0.147} = \qty{706}{m}$ en de snelheid
\qty{10.4}{mm.yr^{-1}}. De twee impliceren zeer verschillende resterende indeukingen.
Een ongecompenseerde indeuking is een massatekort van mantelgesteente: over een breed gebied geeft zij een \omterm{def:b1:isostasy-gravity:anomalies}{vrije-luchtanomalie} van ongeveer $-2\pi G\rho_m w$, dat wil zeggen
$-0.0419\times3.3\times706 \approx \qty{-98}{\milli\gal}$ voor het trage geval, tegenover \qty{-4}{\milli\gal} voor het snelle. Een zwaartekrachtmeting (of de hoogte van oude kustlijnen) onderscheidt ze. (d) De tijdconstante meet hoe snel de \omterm{def:b1:layered-earth:layers}{mantel} onder de indeuking terugstroomt: haar viscositeit.
\end{solution}

\begin{solution}{pb:b1:isostasy-gravity:1}
\textbf{1.} $r = h\rho_c/(\rho_m - \rho_c) = 4\times2700/600 = \qty{18}{km}$.

\textbf{2.} De \omterm{def:b1:layered-earth:layers}{korst} is $35 + 4 + 18 = \qty{57}{km}$ dik; haar basis, de \omterm{def:b1:layered-earth:discontinuity}{Moho}, ligt $35 + 18 = \qty{53}{km}$ onder zeeniveau.

\textbf{3.} Plateau: \qty{57}{km} \omterm{def:b1:layered-earth:layers}{korst} en $60 - 53 = \qty{7}{km}$ \omterm{def:b1:layered-earth:layers}{mantel},
$\num{57000}\times2700 + \num{7000}\times3300 = \num{1.539e8} + \num{2.31e7} =
\qty{1.770e8}{kg.m^{-2}}$. Laagland: $\num{35000}\times2700 + \num{25000}\times3300
= \num{9.45e7} + \num{8.25e7} = \qty{1.770e8}{kg.m^{-2}}$. Ze zijn gelijk, zoals \cref{prop:b1:isostasy-gravity:equal-mass} vereist.

\textbf{4.} $p = g\times\text{massa per oppervlakte-eenheid} = 9.81\times\num{1.770e8} =
\qty{1.74e9}{Pa}$, ongeveer \qty{1.7}{GPa}, onder beide kolommen.

\textbf{5.} De wortel moet door haar kleinere dichtheid de massa van het plateau compenseren: elke meter wortel verwijdert $\rho_m - \rho_c = \qty{600}{kg.m^{-3}}$ aan massa, en elke meter plateau voegt \qty{2700}{kg.m^{-3}} toe, dus zijn $2700/600 = 4.5$ meter wortel nodig. Met een \omterm{def:b1:layered-earth:layers}{korst} van \qty{2900}{kg.m^{-3}} zou de verhouding $2900/400 = 7.25$ zijn: hoe kleiner het contrast, hoe dieper de wortel.

\textbf{6.} De \omterm{def:b1:layered-earth:discontinuity}{Moho}, de sterke toename van de seismische snelheid, zou blijken onder de rand van het plateau dieper te worden, van \qty{35}{km} onder het laagland tot ongeveer \qty{53}{km} onder het plateau.

\textbf{7.} $\rho = 2700\times100/104 = \qty{2596}{kg.m^{-3}}$.

\textbf{8.} $\num{104000}\times2596.2 = \qty{2.700e8}{kg.m^{-2}}$ en
$\num{100000}\times2700 = \qty{2.700e8}{kg.m^{-2}}$.

\textbf{9.} $2700 - 2596 = \qty{104}{kg.m^{-3}}$, of $104/2700 =
\qty{3.85}{\percent}$.

\textbf{10.} $\Delta\rho/\rho = \alpha\,\Delta T$ geeft $\Delta T =
0.0385/\num{3e-5} = \qty{1280}{K}$ gemiddeld over de hele \qty{104}{km} dikke kolom. Dat zou het grootste deel van de kolom doen smelten: warmte alleen kan dit plateau niet licht genoeg maken, en het mechanisme van Pratt kan hier niet het belangrijkste zijn.

\textbf{11.} De kolommen van Pratt hebben geen wortel: de \omterm{def:b1:layered-earth:discontinuity}{Moho} blijft op ongeveer \qty{35}{km} onder het plateau, zoals onder het laagland.

\textbf{12.} \qty{52}{km} ligt dicht bij de \qty{53}{km} van Airy en ver van de \qty{35}{km} van Pratt: het plateau is vooral door een korstwortel gecompenseerd.

\textbf{13.} $0.0419\times2.7\times4000 = \qty{452.5}{\milli\gal}$.

\textbf{14.} Ongecompenseerd (\cref{prop:b1:isostasy-gravity:broad}):
$\Delta g_{\mathrm{FA}} \approx \qty{+452.5}{\milli\gal}$,
$\Delta g_{\mathrm B} \approx 0$.

\textbf{15.} Gecompenseerd: $\Delta g_{\mathrm{FA}} \approx 0$ en
$\Delta g_{\mathrm B} \approx \qty{-452.5}{\milli\gal}$. Het ontbrekende massa van de wortel geeft $-0.0419\times0.6\times\num{18000} = \qty{-452.5}{\milli\gal}$, hetzelfde, omdat $\Delta\rho\,r = \rho_c h$.

\textbf{16.} $\sin^2\ang{30} = 0.25$ en $\sin^2\ang{60} = 0.75$:
$\gamma = 9.780327\times(1 + 0.0013256 - 0.0000044) = \qty{9.793249}{m.s^{-2}}$.
Dan is $g - \gamma = 9.781085 - 9.793249 = \qty{-1216.4}{\milli\gal}$; de \omterm{def:b1:isostasy-gravity:anomalies}{vrije-luchtcorrectie} is $0.3086\times4000 = \qty{1234.4}{\milli\gal}$, dus $\Delta g_{\mathrm{FA}} = \qty{+18.0}{\milli\gal}$, en $\Delta g_{\mathrm B} = 18.0 - 452.5 = \qty{-434.5}{\milli\gal}$.

\textbf{17.} $434.5/452.5 = 0.96$: het plateau is ongeveer \qty{96}{\percent} gecompenseerd (de kleine positieve \omterm{def:b1:isostasy-gravity:anomalies}{vrije-luchtanomalie} is de ongecompenseerde \qty{4}{\percent}).

\textbf{18.} Nee. Nabij het midden van een breed plateau werkt elke anomale laag als een plaat, waarvan de aantrekking niet van haar diepte afhangt: alleen het massatekort per oppervlakte-eenheid telt, en dat is in beide modellen hetzelfde, $\rho_c h$. De twee verschillen alleen nabij de randen, waar de diepte van de compenserende massa's de vorm van de anomalie verandert.

\textbf{19.} Ver van de rand aan beide kanten is de \omterm{def:b1:isostasy-gravity:anomalies}{vrije-luchtanomalie} bijna nul. Een paar kilometer binnen de rand is het plateau onder het station dichtbij en strekt het zich maar aan één kant ver uit: het trekt nog bijna met zijn volle plaatwaarde. De wortel, \qty{35}{km} of meer diep, wordt van daar gezien als een lichaam dat bijna recht onder eindigt: hij trekt met weinig meer dan de helft van zijn effect. Het verschil is een positieve piek. Een paar kilometer buiten de rand trekt het plateau heel weinig, terwijl de wortel nog bijna de helft van zijn effect geeft: een negatieve trog. Het profiel is een positief-negatief paar over elke rand, zoals in de figuur van het hoofdstuk.

\textbf{20.} De \omterm{def:b1:layered-earth:layers}{korst} is nu $57 - 1 = \qty{56}{km}$ dik, $\qty{21}{km}$ meer dan dat van het laagland. In evenwicht splitst dit overschot zich als $h' + r'$ met $r' = 4.5\,h'$: $h' = 21/5.5 = \qty{3.818}{km}$ en $r' = \qty{17.18}{km}$.

\textbf{21.} Het plateau staat nu op \qty{3818}{m}.

\textbf{22.} Het gesteente dat nu aan het oppervlak ligt, lag \qty{1}{km} onder het oude oppervlak, op \qty{3000}{m}; het ligt nu op \qty{3818}{m}: het is met \qty{818}{m} omhooggekomen, dat wil zeggen $1000\times\rho_c/\rho_m$.

\textbf{23.} Elke kilometer geërodeerd gesteente verlaagt het oppervlak met $1 - 2700/3300 =
0.182$ kilometer, dus het naar beneden brengen met \qty{4}{km} kost $4/0.182 = \qty{22}{km}$ erosie, het hele overschot aan dikte $h + r$. In een oude, afgesleten bergketen lagen de gesteenten die nu aan het oppervlak liggen dus ooit ongeveer twintig kilometer diep: de diepe wortels van de gordel komen aan het licht.

\textbf{24.} Het oppervlak daalt met $1000 - 818$: voor elke kilometer gesteente die wordt geërodeerd is de netto daling van het oppervlak
\textbf{\qty[text-series-to-math=true]{182}{m}}.
\end{solution}
